% !TeX root = ../../Book1.tex
% 纲要标题：### 19.5 经典紧群的完整算例
% 纲要位置：几何学计划纲要.md，第 874 行。
% 正文状态：初稿；尚未正式审核。

% 纲要内容开始
% 
% * 酉群和特殊酉群的环面与根
% * 正交群与辛群的根模型
% * 低秩同构、中心及基本群
% 
% ---
% 
% 纲要内容结束

\section{经典紧群的完整算例}
\label{g1:sec:1905}

抽象理论把共轭问题集中到一个环面，再把环面内剩余的识别交给 Weyl 群。对矩阵群，这两步分别对应选择正交或酉基，以及重新排列特征方向。本节把环面、根、整数格、中心和基本群放在同一组模型中计算。根描述伴随作用的无穷小旋转，整数格则记录哪些旋转已经回到单位元；两者同时出现，才不会把具有相同 Lie 代数的群误认为同一个群。

\subsection{酉群和特殊酉群的环面与根}
\label{g1:subsec:1905:01}
先看最熟悉的酉对角化。记 \(\mathfrak u(n)=\{X:X^*=-X\}\)，并令 \(E_{jk}\) 为第 \((j,k)\) 项等于一、其余项为零的矩阵。

\begin{proposition}[酉群的环面与根]\label{g1:prop:unitary-torus-roots}
	设 \(n\ge2\)。在 \(\operatorname{U}(n)\) 中，
	\[
		T=\{\operatorname{diag}(e^{\mathrm{i}\theta_1},\ldots,e^{\mathrm{i}\theta_n}):\theta_j\in\mathbb R\}
	\]
	是极大环面，其 Lie 代数用 \(H(\theta)=\mathrm{i}\operatorname{diag}(\theta_1,\ldots,\theta_n)\) 识别为 \(\mathbb R^n\)。若 \(\varepsilon_j(H(\theta))=\theta_j\)，则
	\[
		\Lambda=2\pi\mathbb Z^n,\qquad
		R=\{\varepsilon_j-\varepsilon_k:j\ne k\},\qquad
		W\cong S_n.
	\]
	在 \(\operatorname{SU}(n)\) 中，相应极大环面为 \(T_0=T\cap\operatorname{SU}(n)\)，Lie 代数及整数格分别是
	\[
		\mathfrak t_0=\{H(\theta):\textstyle\sum_j\theta_j=0\},\qquad
		\Lambda_0=\{2\pi(m_1,\ldots,m_n):m_j\in\mathbb Z,\ \textstyle\sum_jm_j=0\}.
	\]
	其根是上述根在 \(\mathfrak t_0\) 上的限制，Weyl 群仍为 \(S_n\)。
\end{proposition}
\begin{proof}
	酉矩阵有酉特征基。这里所需的谱论可以从一个单位特征向量开始归纳：复特征值存在，酉性使其模为一；该向量的正交补在矩阵及其逆下都不变，因而可继续选基。于是每个酉矩阵酉共轭于 \(T\)。改变对角元的角度给出从单位矩阵到它的路径，再作同一共轭，说明 \(\operatorname{U}(n)\) 连通。对于特殊酉矩阵，特征角之和属于 \(2\pi\mathbb Z\)；给其中一个角减去这个整数倍后，可使角度和恰为零，路径 \(\operatorname{diag}(e^{\mathrm{i}t\theta_j})\) 就始终在 \(\operatorname{SU}(n)\) 中。对角化所用的酉矩阵也可右乘一个对角酉矩阵，将其行列式调成一而不改变对角化。因此 \(\operatorname{SU}(n)\) 连通，且每个元素都在该群内共轭入 \(T_0\)。两群由矩阵方程定义为有界闭集，所以紧。

	\(T_0\) 可由前 \(n-1\) 个相位自由参数化，最后一个取其乘积的逆，故是连通环面。在两种情形，与所有这些对角矩阵交换的矩阵必须是对角矩阵：对第 \((j,k)\) 项选取第 \(j,k\) 个相位不同的环面元即可。于是两群中的中心化子分别就是 \(T,T_0\)，任何包含它们的环面也不能更大。这直接验证了极大性。

	\(\mathfrak u(n)\) 的复化是 \(\mathfrak{gl}(n,\mathbb C)\)，而 \(\mathfrak{su}(n)\) 的复化是 \(\mathfrak{sl}(n,\mathbb C)\)。逐项计算得
	\[
		[H(\theta),E_{jk}]=\mathrm{i}(\theta_j-\theta_k)E_{jk}.
	\]
	故每个非对角矩阵单位给出所列根，零权空间是对角子代数；取迹零只删去中心的一维零权方向。指数是逐项取指数，所以整数格正是所写的集合。

	正规化环面的酉矩阵必须置换其共同特征直线 \(\mathbb Ce_j\)，因而是一个置换矩阵与对角酉矩阵的乘积；反之这样的矩阵确实正规化环面。模掉对角部分只剩置换。在特殊酉群中，任意置换矩阵都能乘以对角酉矩阵将行列式调成一，仍实现相同的置换。因此两个商群都是 \(S_n\)。
\end{proof}

上述根系在迹零超平面内通常记作 \(A_{n-1}\)。一个 Weyl 室可取 \(\theta_1>\theta_2>\cdots>\theta_n\)，在特殊酉情形还要求和为零。相邻两角相等的墙记录两个特征方向不能再由伴随作用区分。\(\operatorname{U}(n)\) 还保留所有角同时增加的中心方向，所以它的根不张成整个 \(\mathfrak t^*\)。

\begin{example}\label{g1:exm:su2-so3-lattices}
	在 \(\operatorname{SU}(2)\) 中写 \(H(t)=\operatorname{diag}(\mathrm{i}t,-\mathrm{i}t)\)。正根为 \(\alpha(H(t))=2t\)，而 \(\exp H(t)=I\) 当且仅当 \(t\in2\pi\mathbb Z\)。商群 \(\operatorname{SU}(2)/\{\pm I\}\) 使用同一个 Lie 代数与同一条根，却在 \(t\in\pi\mathbb Z\) 时就回到单位元。因此根空间完全相同，整数格却扩大一倍。前一个环面的特征标为 \(e^{\mathrm{i}mt}\)，\(m\in\mathbb Z\)；能下降到商环面的特征标还必须在 \(t=\pi\) 取值一，即 \(m\) 必须为偶数。这个商群与 \(\operatorname{SO}(3)\) 的同构将在本节后面直接构造。
\end{example}

\subsection{正交群与辛群的根模型}
\label{g1:subsec:1905:02}
正交群的实特征方向未必是一维。非实特征值成共轭对，对应实二维旋转平面；这使它的环面参数成对出现。令
\[
	R(\theta)=\begin{pmatrix}\cos\theta&-\sin\theta\\\sin\theta&\cos\theta\end{pmatrix},\qquad
	J=\begin{pmatrix}0&-1\\1&0\end{pmatrix}.
\]

\begin{proposition}[旋转群的根模型]\label{g1:prop:orthogonal-torus-roots}
	令 \(J=\begin{pmatrix}0&-1\\1&0\end{pmatrix}\)。在 \(\operatorname{SO}(2r)\)（\(r\ge2\)）中取由旋转块 \(\exp(\theta_jJ)\) 组成的分块对角环面，在 \(\operatorname{SO}(2r+1)\)（\(r\ge1\)）中另加一个末尾的 \(1\) 块。以 \(H(\theta)=\operatorname{diag}(\theta_1J,\ldots,\theta_rJ)\) 表示其 Lie 代数，奇数维时另加零块，并置 \(\varepsilon_j(H(\theta))=\theta_j\)。这两个环面都是极大的，整数格均为 \(2\pi\mathbb Z^r\)，其根与 Weyl 群为
	\[
		\begin{array}{c|c|c}
		G&R&W\\ \hline
		\operatorname{SO}(2r+1)&\{\pm\varepsilon_j\pm\varepsilon_k:j<k\}\cup\{\pm\varepsilon_j\}&(\mathbb Z/2)^r\rtimes S_r\\
		\operatorname{SO}(2r)&\{\pm\varepsilon_j\pm\varepsilon_k:j<k\}&(\mathbb Z/2)^{r-1}\rtimes S_r.
		\end{array}
	\]
	偶数维一行中的符号变换要求改变偶数个坐标的符号。
\end{proposition}
\begin{proof}
	\cref{g1:thm:orthogonal-blocks}的实正交标准分块给出旋转块及 \(1,-1\) 块。行列式为一迫使 \(-1\) 块的数目为偶数，可将它们两两写成 \(R(\pi)\)；\(1\) 块也两两写成 \(R(0)\)，奇数维只留一个 \(1\)。若对角化基的定向不合要求，反转其中一根基向量即可，至多把一个旋转角变号。于是可在特殊正交群内共轭入所列环面。沿各角的线性路径回到零，说明特殊正交群连通；它由正交方程定义为有界闭集，所以紧。

	为计算伴随作用，把标准实表示复化。第 \(j\) 个旋转平面有权 \(\varepsilon_j,-\varepsilon_j\)，奇数维还有一个零权方向。记复双线性标准配对为 \(B\)。映射
	\[
		u\wedge v\longmapsto\bigl(w\longmapsto B(u,w)v-B(v,w)u\bigr)
	\]
	给出 \(\Lambda^2\mathbb C^m\cong\mathfrak{so}(m,\mathbb C)\)，其中 \(m=2r\) 或 \(2r+1\)。在标准基上，它把 \(e_j\wedge e_k\) 送到基本反对称矩阵，因此是线性同构；保持 \(B\) 的矩阵与此映射交换，故还是表示的同构。外积中两权相加，于是不同旋转平面贡献 \(\pm\varepsilon_j\pm\varepsilon_k\)，奇数维的零权方向再贡献 \(\pm\varepsilon_j\)。同一平面的正负权相加为零，共有 \(r\) 个零权方向，恰好构成 \(\mathfrak t_{\mathbb C}\)。因此与 \(\mathfrak t\) 交换的 Lie 代数就是 \(\mathfrak t\)，任何更大的环面都会给出更大的交换子代数，矛盾。

	正规化子置换这些实旋转平面，并可在各平面上反转定向，分别对应置换 \(\theta_j\) 及改变其符号。置换两个二维平面的行列式为一，反转一个平面的行列式为负。在奇数维，可同时反转剩余的一维轴以恢复总行列式一，所以任意符号变化都可实现；偶数维没有这根轴，必须反转偶数个平面。对平面排列及其定向作用平凡的部分就是各平面内的旋转，模掉环面后不再有其他元素。这证明了两个 Weyl 群公式。
\end{proof}

奇数维与偶数维的上述根模型分别记作 \(B_r,D_r\)。\(\operatorname{SO}(2)\cong S^1\) 则是无根的交换群，不能把它当作有非零根的普通 \(D_r\) 例子。\(D_2\) 的四个根分成两组互相正交的正负根，这将在 \(\operatorname{SO}(4)\) 的双覆盖中得到解释。

\ProseParagraphBreak

本节的\term[jinxinqun]{紧辛群}[compact symplectic group]用 \(\operatorname{Sp}(r)\) 表示，它是保持四元数 Hermitian 内积的群，与保持实反对称双线性型的非紧群 \(\operatorname{Sp}(2r,\mathbb R)\) 不同。四元数代数 \(\mathbb H\) 有实基 \(1,\mathbf i,\mathbf j,\mathbf k\)，满足 \(\mathbf i^2=\mathbf j^2=\mathbf k^2=-1\)、\(\mathbf i\mathbf j=\mathbf k=-\mathbf j\mathbf i\)。共轭把虚部取负，且 \(\overline{ab}=\bar b\bar a\)、\(a\bar a=|a|^2\)，故非零元的逆是 \(\bar a/|a|^2\)。把向量写成右 \(\mathbb H\) 向量空间中的列，定义
\[
	\operatorname{Sp}(r)=\{A\in\operatorname{Mat}_{r\times r}(\mathbb H):A^*A=I\},\qquad A^*=\overline A^{\mathsf T}.
\]

\begin{proposition}[紧辛群的根模型]\label{g1:prop:symplectic-torus-roots}
	设 \(r\ge1\)，\(\mathbb H\) 为四元数体，固定四元数单位 \(\mathbf i\)，并令 \(\operatorname{Sp}(r)=\{A\in\operatorname{Mat}_{r\times r}(\mathbb H):A^*A=I\}\)，其中 \(A^*\) 表示四元数共轭转置。则 \(\operatorname{Sp}(r)\) 是实维数 \(r(2r+1)\) 的紧连通 Lie 群。其极大环面可取 \(\operatorname{diag}(e^{\mathbf i\theta_1},\ldots,e^{\mathbf i\theta_r})\)，整数格为 \(2\pi\mathbb Z^r\)。令 \(H(\theta)=\operatorname{diag}(\mathbf i\theta_1,\ldots,\mathbf i\theta_r)\)，用 \(\varepsilon_j(H(\theta))=\theta_j\) 表示坐标权，则
	\[
		R=\{\pm\varepsilon_j\pm\varepsilon_k:j<k\}\cup\{\pm2\varepsilon_j\},\qquad
		W\cong(\mathbb Z/2)^r\rtimes S_r.
	\]
\end{proposition}
\begin{proof}
	映射 \(A\mapsto A^*A\) 的值在四元数 Hermitian 矩阵空间中。在 \(A^*A=I\) 处，微分是 \(V\mapsto A^*V+V^*A\)；对任意 Hermitian 矩阵 \(S\)，取 \(V=AS/2\) 就得到 \(S\)。因而正则值定理给出嵌入子流形结构。Hermitian 矩阵有 \(r\) 个实对角元及 \(r(r-1)/2\) 个任意四元数上三角元，实维数为 \(2r^2-r\)，故群的维数为 \(4r^2-(2r^2-r)=r(2r+1)\)。乘法和求逆光滑，紧性来自闭性及每列长度为一。

	四元数 Gram--Schmidt 与实、复情形采用同一投影步骤：从 \(v\) 中减去 \(\sum_j u_j\langle u_j,v\rangle\)，再除以正实范数。右乘系数的次序不能交换。这说明群传递地作用于 \(S^{4r-1}\)，第一列为指定单位向量的稳定子同构于 \(\operatorname{Sp}(r-1)\)。\cref{g1:thm:homogeneous-space,g1:prop:transitive-action-model}把此球面识别为稳定子的商，并给出局部截面，故取第一列是局部平凡丛。\(\operatorname{Sp}(1)=S^3\) 连通；归纳假定纤维连通，沿球面路径用有限个局部平凡化提升，再在终点纤维内连接，就能把任一矩阵连到单位矩阵。因此各 \(\operatorname{Sp}(r)\) 都连通。

	写四元数矩阵为 \(Z+W\mathbf j\)，其中 \(Z,W\) 是复矩阵。代入 \(\mathbf jz=\bar z\mathbf j\) 可验证
	\[
		Z+W\mathbf j\longmapsto
		\begin{pmatrix}Z&W\\-\bar W&\bar Z\end{pmatrix}
	\]
	保持乘法及共轭转置。故 Lie 代数可写成
	\[
		\begin{pmatrix}A&B\\-\bar B&\bar A\end{pmatrix},\qquad A^*=-A,\quad B^{\mathsf T}=B.
	\]
	环面在这个复矩阵模型中为 \(\operatorname{diag}(e^{\mathrm{i}\theta_1},\ldots,e^{\mathrm{i}\theta_r},e^{-\mathrm{i}\theta_1},\ldots,e^{-\mathrm{i}\theta_r})\)。\(A\) 的非对角项给出 \(\varepsilon_j-\varepsilon_k\)，\(B\) 的非对角对称项给出 \(\varepsilon_j+\varepsilon_k\)，其对角项给出 \(2\varepsilon_j\)；共轭块给出各负根。零权空间只有 \(A\) 的纯虚对角元，故所列环面极大。

	置换四元数坐标实现任意 \(S_r\) 作用，在某一坐标左乘单位四元数 \(\mathbf j\)，则由 \(\mathbf j e^{\mathbf i\theta}\mathbf j^{-1}=e^{-\mathbf i\theta}\) 实现单个符号反转。这些矩阵属于紧辛群。反过来，正规化子在标准复表示中只能置换成对的权线 \(\varepsilon_j,-\varepsilon_j\)，所以作用必是带符号的置换；作用平凡的部分是环面。这就确定了 Weyl 群。
\end{proof}

这组根记作 \(C_r\)。\(B_r\) 与 \(C_r\) 有相同的反射超平面与 Weyl 群，但轴向根分别是 \(\varepsilon_j\) 和 \(2\varepsilon_j\)，它们与斜向根的相对长度不同。只画反射墙会丢掉这一差别；比较根模型还必须保留根向量与所用坐标尺度。

\subsection{低秩同构、中心及基本群}
\label{g1:subsec:1905:03}
根与 Weyl 群已经能控制许多共轭计算，但它们尚未决定中心与覆盖关系。我们先构造低维覆盖，再用球面上的标架丛计算基本群，避免把覆盖次数从根图中猜出来。

\begin{proposition}[四元数给出的低维群映射]\label{g1:prop:quaternion-low-rank}
	有光滑群同构 \(\operatorname{Sp}(1)\cong\operatorname{SU}(2)\)，其共同底流形为 \(S^3\)。另有满射 Lie 群同态
	\[
		\operatorname{Sp}(1)\longrightarrow\operatorname{SO}(3),\qquad
		\operatorname{Sp}(1)\times\operatorname{Sp}(1)\longrightarrow\operatorname{SO}(4),\qquad
		\operatorname{Sp}(2)\longrightarrow\operatorname{SO}(5),
	\]
	其核依次为 \(\{\pm1\}\)、\(\{(1,1),(-1,-1)\}\)、\(\{\pm I\}\)，所以均为二重覆盖。
\end{proposition}
\begin{proof}
	单位四元数 \(z+w\mathbf j\) 对应
	\[
		\begin{pmatrix}z&w\\-\bar w&\bar z\end{pmatrix},\qquad |z|^2+|w|^2=1.
	\]
	其列正交且行列式为一；反过来，特殊酉二阶矩阵的第二行由第一行及行列式一唯一确定为这一形式。上面已经核对乘法相容，所以得到光滑群同构及 \(S^3\) 模型。

	单位四元数在三维虚四元数空间上作用为 \(v\mapsto qvq^{-1}\)。共轭保持虚部与范数，且从单位元连续变化，故像在 \(\operatorname{SO}(3)\)。若作用恒等，\(q\) 与 \(\mathbf i,\mathbf j,\mathbf k\) 都交换，逐项比较便知 \(q\) 是实数，故核为 \(\{\pm1\}\)。源、靶同为三维，离散核使单位元微分单射，因而为同构。像是开子群，又因源紧而闭；靶连通，所以像为整个群。\cref{g1:prop:lie-homomorphism-local-isomorphism}于是给出二重覆盖。也可直接令 \(q=\cos t+u\sin t\)，其中 \(u\) 是单位虚四元数：共轭固定 \(u\)，在 \(u^\perp\) 上旋转 \(2t\)。这解释了\cref{g1:exm:su2-so3-lattices}中的两倍角。

	在四维实空间 \(\mathbb H\) 上令 \((a,b)\) 作用为 \(x\mapsto axb^{-1}\)。它保持范数，且源连通，像在 \(\operatorname{SO}(4)\)。若作用恒等，代入 \(x=1\) 得 \(a=b\)，再代入所有四元数得 \(a=b=\pm1\)。两群维数均为六，刚才的开闭像论证再次适用，得到第二个覆盖。

	最后令 \(V\) 为迹零的二阶四元数 Hermitian 矩阵空间：
	\[
		V=\left\{\begin{pmatrix}r&q\\\bar q&-r\end{pmatrix}:r\in\mathbb R,\ q\in\mathbb H\right\},\qquad
		\|X\|^2=\tfrac12\operatorname{Re}\operatorname{tr}(X^2)=r^2+|q|^2.
	\]
	它是五维实 Euclid 空间。\(A\in\operatorname{Sp}(2)\) 的作用为 \(X\mapsto AXA^*\)。四元数满足 \(\operatorname{Re}(ab)=\operatorname{Re}(ba)\)，按分量求和便得到实部迹的循环性；因而这个作用保持迹零条件及所写范数。源连通，像落在 \(\operatorname{SO}(V)\)。核中的 \(A\) 必须与 \(\operatorname{diag}(1,-1)\) 交换，所以 \(A=\operatorname{diag}(a,d)\)；再与所有非对角 \(X\) 交换，得 \(aq=qd\) 对所有 \(q\in\mathbb H\) 成立。取 \(q=1,\mathbf i,\mathbf j\) 得 \(a=d=\pm1\)。两群维数都是十，同样得到满射二重覆盖。
\end{proof}

\begin{lemma}[球面上的标架丛与基本群]\label{g1:lem:sphere-bundle-pi1}
	设 \(K\to E\xrightarrow{p}S^m\) 是光滑主 \(K\) 丛，\(K\) 为连通 Lie 群，\(m\ge3\)。则纤维包含 \(K\hookrightarrow E\) 诱导基本群同构。
\end{lemma}
\begin{proof}
	取分别去掉南、北极的开集 \(U,V\)，二者都微分同胚于 \(\mathbb R^m\)，交集同伦等价于 \(S^{m-1}\)。丛在 \(U,V\) 上各有全局截面。为说明这一点，可按\cref{g1:thm:principal-connection-existence}取主联络，在 Euclid 坐标中从原点向每点作径向路径，再从固定纤维点水平提升；\cref{g1:thm:principal-horizontal-lift}保证提升在整个区间存在并光滑依赖终点，因此给出所需截面。这只用主联络沿紧区间的提升，不要求底空间完备。

	于是 \(p^{-1}(U)\) 与 \(p^{-1}(V)\) 都是可缩开集与 \(K\) 的乘积；交集用任一平凡化写成 \((U\cap V)\times K\)。因 \(m-1\ge2\)，\cref{g1:prop:sphere-simply-connected}说明 \(U\cap V\) 单连通。乘积中的路径与同伦逐分量进行，所以这三个空间的基本群都通过同一基点纤维与 \(\pi_1(K)\) 同构。两条包含映射在该纤维上就是恒等，故在这些识别下也都是恒等。\cref{g1:thm:van-kampen}将两个 \(\pi_1(K)\) 沿共同的同一 \(\pi_1(K)\) 粘合，所得群仍为 \(\pi_1(K)\)，且同构正由纤维包含诱导。
\end{proof}

\begin{proposition}[经典紧群的中心与基本群]\label{g1:prop:classical-centers-pi1}
	下表列出相应连通经典紧群的中心与基本群。其中 \(\mu_n=\{z\in S^1:z^n=1\}\)，\(r\ge1\)，并单独保留低维旋转群。
	\[
		\begin{array}{c|c|c}
		G&Z(G)&\pi_1(G)\\ \hline
		\operatorname{U}(n),\ n\ge1&S^1I&\mathbb Z\\
		\operatorname{SU}(n),\ n\ge2&\mu_nI&0\\
		\operatorname{Sp}(r)&\{\pm I\}&0\\
		\operatorname{SO}(2)&\operatorname{SO}(2)&\mathbb Z\\
		\operatorname{SO}(2r+1),\ r\ge1&\{I\}&\mathbb Z/2\\
		\operatorname{SO}(2r),\ r\ge2&\{\pm I\}&\mathbb Z/2.
		\end{array}
	\]
\end{proposition}
\begin{proof}
	酉或特殊酉群的中心元必须与其对角环面交换，所以为对角矩阵；再与混合任意两个坐标的特殊酉二阶块交换，迫使所有对角元相同。行列式条件分别给出 \(S^1I\) 与 \(\mu_nI\)。紧辛群中，先与复对角环面交换得到复对角矩阵，再与各坐标上的 \(\mathbf j\) 交换，迫使对角元为实数；置换坐标使它们相同，故只有 \(\pm I\)。对于 \(m\ge3\)，若实线性变换与所有坐标平面旋转交换，则对每个单位向量 \(v\)，它必须保持固定 \(v\) 的旋转子群的共同不动空间 \(\mathbb Rv\)。因此每条直线不变，由线性性是标量变换。正交性及行列式一分别给出奇数维的 \(I\) 和偶数维的 \(\pm I\)。\(\operatorname{SO}(2)\) 本来交换，须另行处理。

	取第一列给出主丛
	\[
		\operatorname{SU}(n-1)\longrightarrow\operatorname{SU}(n)\longrightarrow S^{2n-1},\qquad
		\operatorname{Sp}(r-1)\longrightarrow\operatorname{Sp}(r)\longrightarrow S^{4r-1}.
	\]
	第一条中的传递性由酉 Gram--Schmidt 得到：选定第一列后，用其余列之一调整行列式而不改变第一列。稳定子恰为低一阶的特殊酉群。第二条的传递性和稳定子已经在\cref{g1:prop:symplectic-torus-roots}中证明。它们均由齐性空间的局部截面构成主丛，符合\cref{g1:exm:homogeneous-principal-bundle}的条件。以 \(\operatorname{SU}(1)=\operatorname{Sp}(0)=\{1\}\) 为起点，反复应用上一引理，得到两类群的基本群都为零。\(\operatorname{SU}(2)\cong S^3\) 还给出独立的起始核验。

	映射
	\[
		S^1\times\operatorname{SU}(n)\longrightarrow\operatorname{U}(n),\qquad
		(z,A)\longmapsto\operatorname{diag}(z,1,\ldots,1)A
	\]
	是微分同胚：逆映射先取 \(z=\det U\)，再取 \(A=\operatorname{diag}(z^{-1},1,\ldots,1)U\)。它通常不是直积群同构，但已足够说明 \(\pi_1(\operatorname{U}(n))\cong\pi_1(S^1)=\mathbb Z\)，且行列式诱导此同构。

	\(S^3\to\operatorname{SO}(3)\) 是前一命题给出的二重覆盖；\(S^3\) 单连通，所以\cref{g1:thm:normal-covering-deck}将基本群识别为二元覆盖变换群，给出 \(\pi_1(\operatorname{SO}(3))\cong\mathbb Z/2\)。对 \(n\ge4\)，第一列又给出主丛
	\[
		\operatorname{SO}(n-1)\longrightarrow\operatorname{SO}(n)\longrightarrow S^{n-1}.
	\]
	正交标架可在不改变第一列的条件下调整定向，故作用传递，稳定子为 \(\operatorname{SO}(n-1)\)。球面维数至少三，上一引理说明包含映射诱导基本群同构，归纳便得全部 \(n\ge3\) 的结论。\(\operatorname{SO}(2)\) 的一整周旋转生成 \(\mathbb Z\)，进入 \(\operatorname{SO}(3)\) 后只保留其奇偶性。
\end{proof}

至此，\(\operatorname{Sp}(1)\times\operatorname{Sp}(1)\to\operatorname{SO}(4)\) 与 \(\operatorname{Sp}(2)\to\operatorname{SO}(5)\) 的源群也已证明单连通，因而是相应的万有覆盖。微分给出 \(\mathfrak{so}(4)\cong\mathfrak{su}(2)\oplus\mathfrak{su}(2)\) 及 \(\mathfrak{so}(5)\cong\mathfrak{sp}(2)\)，但群本身仍须除以所列二元中心核。根模型中的 \(D_2\) 分裂及 \(B_2,C_2\) 的低秩对应因此各有具体的矩阵实现。

\ProseParagraphBreak

这些算例显示三种资料各自回答不同问题。根给出共轭轨道的无穷小方向，Weyl 群给出同一环面内的有限识别，整数格与中心商则决定哪些环路真正闭合以及覆盖会怎样改变群。后一类信息不能从 Lie 代数同构中自动恢复。
